\input{configTD}
\usepackage{tikz}
\usepackage{needspace}

% ── Poly à trous : commandes spécifiques ──────────────────────────
\newcommand{\atrou}[2][3cm]{%
  \ifthenelse{\boolean{showSolutions}}{%
    \par\vspace{0.5em}%
    \begin{mdframed}[backgroundcolor=ThemeLight, linewidth=0pt,
      skipabove=0pt, skipbelow=0pt,
      innertopmargin=8pt, innerbottommargin=8pt]%
    #2%
    \end{mdframed}%
    \par\vspace{0.5em}%
  }{%
    \par\vspace{0.5em}%
    \noindent\fbox{\begin{minipage}[c][#1]%
      {\dimexpr\textwidth-2\fboxsep-2\fboxrule\relax}%
      \color{white}\rule{0pt}{0pt}%
    \end{minipage}}%
    \par\vspace{0.5em}%
  }%
}

\newcommand{\val}[1]{%
  \ifthenelse{\boolean{showSolutions}}{#1}{\phantom{#1}}%
}

\newcommand{\R}{\mathbb{R}}
\newcommand{\E}{\mathbb{E}}
\newcommand{\Var}{\operatorname{Var}}
\newcommand{\Xbar}{\overline{X}_n}

\title{\sffamily\bfseries Probabilités --- CM5\\[0.3em]
  \Large Tests d'hypothèses : à partir de quand être surpris\,?}
\author{}
\date{}

\begin{document}
\sffamily
\maketitle
\thispagestyle{empty}

% ====================================================================
\section*{1. Un problème : faut-il être surpris\,?}
% ====================================================================

\textbf{Cadre.} Vous êtes au volant sur l'autoroute. Vous savez (statistiques constructeurs) qu'en France, environ $10\,\%$ des voitures sont \textbf{rouges}. Soudain, vous croisez \textbf{$4$ voitures rouges de suite}.

\medskip

\begin{center}
\begin{mdframed}[backgroundcolor=ThemeLight, linecolor=Theme,
  linewidth=1.5pt, innertopmargin=10pt, innerbottommargin=10pt]
\centering\large\itshape
À partir de quand commence-t-on à douter du modèle\\
\og les couleurs des voitures qui passent sont indépendantes,\\
chacune rouge avec probabilité $0{,}10$\,\fg{}\,?
\end{mdframed}
\end{center}

\textbf{Question 1.} Sous ce modèle, quelle est la probabilité de croiser $4$ voitures rouges \emph{exactement} d'affilée\,?

\atrou[2.5cm]{%
Les couleurs étant indépendantes :
\[
  P(\text{4 rouges de suite}) = 0{,}10^4 = 10^{-4} = 0{,}0001.
\]
C'est environ une chance sur $10\,000$. \emph{Rare}, mais pas impossible.
}

\textbf{Question 2.} Et pour $7$ voitures rouges d'affilée\,?

\atrou[2cm]{%
\[
  P(\text{7 rouges de suite}) = 0{,}10^7 = 10^{-7} = 0{,}0000001.
\]
Une chance sur dix millions. À ce stade, on ne croit \emph{plus} au modèle.
}

\medskip

\begin{mdframed}[backgroundcolor=ThemeLight, linecolor=Theme,
  linewidth=1.5pt, innertopmargin=10pt, innerbottommargin=10pt]
\textbf{La logique du test d'hypothèse --- en quatre temps.}
\begin{enumerate}
  \item \textbf{Hypothèse de référence} ($H_0$). On pose un modèle \og par défaut \fg{} dont on veut éprouver la crédibilité (ici : indépendance + $p=0{,}10$).
  \item \textbf{Observation}. On collecte des données (ici : $4$ rouges de suite).
  \item \textbf{Probabilité sous $H_0$}. On calcule la probabilité, \emph{en supposant $H_0$ vraie}, d'observer un résultat \og au moins aussi extrême \fg{} que celui constaté. C'est la \textbf{p-value}.
  \item \textbf{Décision}. Si cette probabilité est \emph{très petite} (sous un seuil $\alpha$ fixé à l'avance), on \textbf{rejette $H_0$} : les données sont incompatibles avec le modèle. Sinon, on \emph{conserve} $H_0$ par défaut.
\end{enumerate}
\end{mdframed}

\medskip

\textit{Remarque.} On ne \og prouve \fg{} jamais qu'une hypothèse est vraie. On peut seulement décider qu'elle est trop incompatible avec les données pour être maintenue. Comme au tribunal : on rejette la présomption d'innocence (au-delà d'un doute raisonnable) ou bien on l'\textbf{accepte par défaut} --- ce n'est pas pareil que \og prouver que l'accusé est innocent \fg.

\subsection*{1.1. Deuxième exemple : un dé sans aucun \texorpdfstring{$6$}{6}}

On lance un dé à $6$ faces $20$ fois de suite et, à notre grande surprise, on n'obtient \textbf{aucun $6$}. Doit-on suspecter que le dé est truqué\,?

\textbf{Question 1.} Énoncer l'hypothèse $H_0$ (\og le dé n'est pas truqué \fg) et proposer une statistique de test naturelle. Quelle est sa loi sous $H_0$\,?

\atrou[3.5cm]{%
$H_0$ : \og le dé est équilibré \fg, c'est-à-dire $P(\text{face } 6) = 1/6$ et lancers indépendants.\\
Statistique naturelle : $X$ = nombre de $6$ obtenus en $20$ lancers. Sous $H_0$,
\[
  X \sim \mathcal{B}\!\left(20,\,\tfrac{1}{6}\right),
  \qquad \E(X) = \tfrac{20}{6} \approx 3{,}33.
\]
}

\textbf{Question 2.} On observe $X_{\text{obs}} = 0$. Calculer la p-value et conclure au seuil $\alpha = 5\,\%$.

\atrou[3.5cm]{%
On observe une valeur \emph{anormalement basse}~; la p-value est la probabilité d'observer aussi peu de $6$ que cela sous $H_0$ :
\[
  \text{p-value} = P(X \le 0 \mid H_0) = P(X = 0) = \left(\tfrac{5}{6}\right)^{20} \approx 0{,}026.
\]
$0{,}026 < 0{,}05$ : on \textbf{rejette} $H_0$. Le dé semble truqué (manque de $6$).

\smallskip

\emph{Anecdote.} Au seuil $\alpha = 1\,\%$, on ne rejetterait \emph{pas} : la conclusion dépend du seuil choisi --- d'où l'importance de le fixer avant l'expérience.
}

% ====================================================================
\section*{2. Formalisation : hypothèse, p-value, décision}
% ====================================================================

\begin{mdframed}[backgroundcolor=ThemeLight, linecolor=Theme,
  linewidth=1pt, innertopmargin=8pt, innerbottommargin=8pt, nobreak=true]
\textbf{Vocabulaire.}
\begin{itemize}
  \item \textbf{Hypothèse nulle} $H_0$ : l'hypothèse \og par défaut \fg{} qu'on cherche à mettre à l'épreuve.
  \item \textbf{Hypothèse alternative} $H_1$ : ce qui reste si $H_0$ est fausse.
  \item \textbf{Statistique de test} : une quantité calculée sur l'échantillon, dont on connaît la loi \emph{sous} $H_0$.
  \item \textbf{p-value} : probabilité, \emph{sous $H_0$}, d'observer un résultat au moins aussi extrême que celui qu'on a obtenu.
  \item \textbf{Seuil} $\alpha$ : seuil de tolérance (typiquement $5\,\%$ ou $1\,\%$). On rejette $H_0$ si p-value $<\alpha$.
\end{itemize}
\end{mdframed}

\medskip

\textbf{Question 1.} Dans l'exemple des voitures, quelles sont $H_0$ et $H_1$\,? Quelle p-value obtient-on avec $4$ rouges d'affilée\,? Au seuil $\alpha = 0{,}05$, rejette-t-on $H_0$\,?

\atrou[3.5cm]{%
$H_0$ : \og les couleurs sont indépendantes et $p_{\text{rouge}} = 0{,}10$ \fg{}.\\
$H_1$ : \og il y a un biais (dépendance ou $p$ plus grand) \fg{}.\\
p-value $= P(\text{au moins 4 rouges de suite} \mid H_0) = 10^{-4}$.\\
$10^{-4} < 0{,}05$, donc on \textbf{rejette} $H_0$.
}

\medskip

\textbf{Le test d'hypothèse, c'est partout.} Le même squelette \og hypothèse de référence \textrightarrow{} statistique \textrightarrow{} p-value \textrightarrow{} décision \fg{} se retrouve dans une foule de domaines très différents. Pour chacun des contextes ci-dessous, à toi de proposer l'hypothèse $H_0$.

\smallskip

\textbf{Exemple 1 --- Médecine.} On veut tester un nouveau traitement contre l'hypertension. On recrute deux groupes de $200$ patients : l'un reçoit le médicament, l'autre un placebo. Au bout de $3$ mois, la baisse moyenne de tension est de $9$~mmHg dans le premier groupe contre $7$~mmHg dans le second. Cette différence est-elle réelle, ou s'explique-t-elle par le simple hasard de l'échantillonnage\,?

\atrou[2cm]{%
$H_0$ : \og le traitement n'a aucun effet \fg{}, soit $\mu_{\text{traitement}} = \mu_{\text{placebo}}$. La différence de $2$~mmHg n'est qu'une fluctuation due aux échantillons.
}

\textbf{Exemple 2 --- Contrôle qualité.} Une chaîne de production garantit au plus $1\,\%$ de boulons défectueux. Un client effectue une réception~: il prélève $500$ boulons et en trouve $11$ défectueux ($2{,}2\,\%$). Le lot doit-il être refusé\,?

\atrou[2cm]{%
$H_0$ : \og la chaîne respecte sa spécification \fg, soit $p_{\text{défaut}} \le 0{,}01$. Le test met à l'épreuve l'hypothèse que la fréquence observée ($2{,}2\,\%$) est compatible avec un vrai $p$ inférieur à $1\,\%$.
}

\textbf{Exemple 3 --- A/B testing.} Le service marketing d'un site marchand teste deux versions de la page d'accueil. La version~A est montrée à $10\,000$ visiteurs ($320$ achats), la version~B à $10\,000$ autres ($360$ achats). Faut-il déployer la version~B partout\,?

\atrou[2cm]{%
$H_0$ : \og les deux versions ont le même taux de conversion \fg, soit $p_A = p_B$. La différence observée ($3{,}2\,\%$ vs $3{,}6\,\%$) pourrait n'être qu'un effet du hasard.
}

\textbf{Exemple 4 --- Sondage électoral.} À $48\,\%$ d'intentions de vote sur un échantillon de $1000$ personnes, le candidat peut-il être assuré qu'il franchira la barre des $50\,\%$ au second tour\,?

\atrou[2cm]{%
$H_0$ : \og le candidat est sous la barre \fg, soit $p \le 0{,}50$. Tant qu'on ne rejette pas $H_0$, on ne peut pas conclure qu'il est majoritaire dans la population.
}

\textbf{Exemple 5 --- Détection de fraude.} Dans une comptabilité \og naturelle \fg, le \emph{premier chiffre} d'un montant suit la \textbf{loi de Benford} ($1 \to 30\,\%$, $2 \to 18\,\%$, \ldots, $9 \to 4{,}5\,\%$). Un enquêteur compare la distribution des premiers chiffres d'une déclaration suspecte à cette loi.

\atrou[2cm]{%
$H_0$ : \og les premiers chiffres suivent la loi de Benford \fg{} (donc pas de manipulation). Un écart \og trop grand \fg{} fait rejeter l'hypothèse, et plaide pour des chiffres fabriqués (les humains tirent uniformément quand ils inventent).
}

\medskip

\begin{mdframed}[backgroundcolor=ThemeLight, linecolor=Theme,
  linewidth=1pt, innertopmargin=8pt, innerbottommargin=8pt]
\textbf{Quelques seuils $\alpha$ \og publics \fg{}.} Le seuil dépend des conséquences d'une fausse découverte~:
\begin{itemize}\setlength{\itemsep}{0.15em}
  \item \textbf{Sciences sociales, médecine usuelle}~: $\alpha = 5\,\%$ (par convention).
  \item \textbf{Essais cliniques de phase~III (FDA, EMA)}~: typiquement $\alpha = 5\,\%$, mais on demande \emph{deux études indépendantes positives} avant de mettre un médicament sur le marché.
  \item \textbf{Études génomiques (GWAS)}~: $\alpha \approx 5 \cdot 10^{-8}$, pour corriger les ${\sim}10^6$ tests menés simultanément sur le génome.
  \item \textbf{Physique des particules (CERN)}~: pour annoncer une \og découverte \fg{} (boson de Higgs, $2012$), il faut $5\sigma$, soit p-value $< 3 \cdot 10^{-7}$.
\end{itemize}
Plus l'enjeu est lourd (autoriser un médicament, annoncer une particule fondamentale), plus le seuil est sévère.
\end{mdframed}

\medskip

\begin{mdframed}[backgroundcolor=ThemeLight, linecolor=Accent,
  linewidth=1.5pt, innertopmargin=8pt, innerbottommargin=8pt]
\textbf{Attention au $p$-hacking.} Pour que la p-value ait du sens, le seuil $\alpha$ \emph{et} la statistique de test doivent être fixés \textbf{avant} de regarder les données. Sinon, on tombe dans des pratiques qui fabriquent artificiellement des résultats \og significatifs \fg{} :
\begin{itemize}\setlength{\itemsep}{0.15em}
  \item \textbf{Tester $20$ hypothèses, garder celle qui marche.} Au seuil $5\,\%$, on s'attend à $1$ résultat \og positif \fg{} par pur hasard sur $20$, même quand tout est faux.
  \item \textbf{Stopper la collecte dès que p $<\,0{,}05$.} En continuant à observer juste assez, on finit toujours par franchir le seuil.
  \item \textbf{Découper en sous-groupes après coup.} \og Le médicament ne marche pas en moyenne, mais marche \emph{chez les femmes de plus de $50$ ans} \fg{} : avec assez de découpages, une catégorie sera \og positive \fg{} par hasard.
\end{itemize}
C'est l'une des causes de la \emph{crise de reproductibilité} en sciences (beaucoup de résultats publiés ne se confirment pas en répétant l'expérience).
\end{mdframed}

% ====================================================================
\section*{3. Test sur une moyenne}
% ====================================================================

On revient sur l'exemple central du CM4 : une usine remplit des sacs de ciment de poids nominal $25$~kg, avec un écart-type de remplissage $\sigma = 0{,}5$~kg \emph{connu}. Le chef de production veut savoir si la machine est \textbf{déréglée}.

Il pèse $n=100$ sacs et obtient une moyenne empirique $\Xbar = 24{,}8$~kg.

\medskip

\textbf{Question 1.} Quelles sont $H_0$ et $H_1$ dans ce problème\,?

\atrou[2cm]{%
$H_0$ : $\mu = 25$ (la machine est bien réglée).\\
$H_1$ : $\mu \neq 25$ (la machine est déréglée).
}

\textbf{Question 2.} En utilisant le TCL, construire une \textbf{statistique de test} $Z$ qui suit $\mathcal{N}(0,1)$ sous $H_0$. Calculer la valeur observée $Z_{\text{obs}}$.

\atrou[4cm]{%
Sous $H_0$, les sacs sont i.i.d.\ d'espérance $25$ et d'écart-type $0{,}5$. Par le TCL,
\[
  \Xbar \;\approx\; \mathcal{N}\!\left(25,\;\dfrac{0{,}5^2}{100}\right).
\]
On standardise pour obtenir une loi $\mathcal{N}(0,1)$ :
\[
  Z = \dfrac{\Xbar - 25}{\sigma/\sqrt{n}} \;\approx\; \mathcal{N}(0,1),
  \qquad
  Z_{\text{obs}} = \dfrac{24{,}8 - 25}{0{,}5/\sqrt{100}} = \dfrac{-0{,}2}{0{,}05} = -4.
\]
}

\textbf{Question 3.} Quelle est la p-value\,? Conclure au seuil $\alpha = 5\,\%$.

\atrou[4.5cm]{%
La machine peut se dérégler vers le haut ou vers le bas : on s'intéresse à l'écart \emph{dans les deux sens}.
\[
  \text{p-value} = P\bigl(|Z| \ge |Z_{\text{obs}}| \;\big|\; H_0\bigr)
  = P(|Z| \ge 4) = 2\bigl(1-\Phi(4)\bigr).
\]
Or $1-\Phi(4) \approx 3{,}2\cdot 10^{-5}$, donc p-value $\approx 6{,}4\cdot 10^{-5}$.\\
$6{,}4\cdot 10^{-5} \ll 0{,}05$ : on \textbf{rejette} $H_0$. La machine est déréglée.
}

\begin{center}
\begin{tikzpicture}[scale=0.85]
  % bell curve
  \draw[thick, Theme, domain=-3.4:3.4, smooth, samples=100]
    plot (\x, {2.2*exp(-\x*\x/2)});
  % shaded tails (z = -1.96 and 1.96)
  \fill[ThemeLight, opacity=0.7, domain=-3.4:-1.96, smooth, samples=40]
    plot (\x, {2.2*exp(-\x*\x/2)}) -- (-1.96,0) -- (-3.4,0) -- cycle;
  \fill[ThemeLight, opacity=0.7, domain=1.96:3.4, smooth, samples=40]
    plot (\x, {2.2*exp(-\x*\x/2)}) -- (3.4,0) -- (1.96,0) -- cycle;
  \draw[thick,->] (-4,0) -- (4,0) node[right] {\small $z$};
  \draw[dashed] (-1.96,0) -- (-1.96,1);
  \draw[dashed] (1.96,0) -- (1.96,1);
  \node at (-1.96,-0.35) {\small $-1{,}96$};
  \node at (1.96,-0.35) {\small $1{,}96$};
  \node at (0,-0.35) {\small $0$};
  \node at (-2.8,0.4) {\small\bfseries $\alpha/2$};
  \node at (2.8,0.4) {\small\bfseries $\alpha/2$};
  \node at (0,2.6) {\small zone d'\textbf{acceptation} (on garde $H_0$)};
\end{tikzpicture}
\end{center}

\medskip

\begin{mdframed}[backgroundcolor=ThemeLight, linecolor=Theme,
  linewidth=1pt, innertopmargin=8pt, innerbottommargin=8pt, nobreak=true]
\textbf{Recette du test sur une moyenne ($\sigma$ connu).}
\begin{enumerate}
  \item $H_0 : \mu = \mu_0$ \quad contre \quad $H_1 : \mu \neq \mu_0$.
  \item Statistique : $Z = \dfrac{\Xbar - \mu_0}{\sigma/\sqrt{n}}$, qui suit $\mathcal{N}(0,1)$ sous $H_0$.
  \item Seuil $\alpha$ (souvent $0{,}05$) $\Longrightarrow$ valeur critique $t_\alpha$ (par exemple $t_{0{,}05}=1{,}96$).
  \item \textbf{Décision} : \textbf{rejeter} $H_0$ si $|Z_{\text{obs}}| > t_\alpha$, sinon \textbf{ne pas rejeter}.
\end{enumerate}
\end{mdframed}

\subsection*{3.1. Exercice : ce dé est-il pipé\,? (test du \texorpdfstring{$\chi^2$}{khi-deux})}

\textbf{Au-delà de la moyenne.} Le test précédent compare une moyenne empirique à une valeur attendue. On veut souvent tester quelque chose de plus riche : \og les fréquences observées d'un phénomène réparti en plusieurs catégories sont-elles compatibles avec celles attendues sous un modèle de référence\,? \fg{} C'est le rôle du \textbf{test du khi-deux d'adéquation} (Karl Pearson, $1900$).

\smallskip

\textbf{Quelques applications classiques.}
\begin{itemize}\setlength{\itemsep}{0.15em}
  \item \textbf{Adéquation à une loi de référence}~: un dé est-il équilibré\,? les premiers chiffres d'une comptabilité suivent-ils la loi de Benford\,?
  \item \textbf{Indépendance entre deux variables qualitatives}~: la couleur des yeux dépend-elle du sexe\,? le secteur d'activité dépend-il de la région\,?
  \item \textbf{Homogénéité entre populations}~: la répartition des opinions est-elle la même chez les $\le 30$~ans et les $\ge 60$~ans\,?
\end{itemize}

\medskip

\begin{mdframed}[backgroundcolor=ThemeLight, linecolor=Theme,
  linewidth=1pt, innertopmargin=8pt, innerbottommargin=8pt]
\textbf{La loi du khi-deux $\chi^2(k)$.} Si $U_1, \dots, U_k$ sont $k$ variables \emph{indépendantes} de loi $\mathcal{N}(0,1)$, alors par définition
\[
  U_1^2 + U_2^2 + \dots + U_k^2 \;\sim\; \chi^2(k).
\]
On l'appelle \og loi du khi-deux à $k$ degrés de liberté \fg{}. Elle ne prend que des valeurs positives, son espérance vaut $k$, son écart-type $\sqrt{2k}$. Plus la valeur observée s'éloigne de $k$ vers la droite, plus elle est \emph{suspecte}.

\begin{center}
\begin{tikzpicture}[xscale=0.4, yscale=12]
  % chi²(5) density: f(x) = (1/(2^{5/2} Gamma(5/2))) * x^{3/2} * exp(-x/2)
  % ≈ 0.0884 * x^{1.5} * exp(-x/2). Mode at x=3, peak ≈ 0.154.
  \draw[thick, Theme, smooth, samples=80, domain=0.05:22]
    plot (\x, {0.0884 * pow(\x,1.5) * exp(-\x/2)});
  \fill[Accent, opacity=0.5, smooth, samples=40, domain=11.07:22]
    plot (\x, {0.0884 * pow(\x,1.5) * exp(-\x/2)})
      -- (22,0) -- (11.07,0) -- cycle;
  \draw[thick,->] (-0.3,0) -- (23,0) node[right] {\scriptsize $x$};
  \draw[dashed, thin] (5,0) -- (5,0.16);
  \draw[dashed, thin] (11.07,0) -- (11.07,0.05);
  \node[below=1pt] at (5,0) {\scriptsize $\E = k = 5$};
  \node[below=1pt] at (11.07,0) {\scriptsize $11{,}07$};
  \node[font=\scriptsize, Accent] at (15.5,0.025) {\bfseries $\alpha = 5\,\%$};
  \node[font=\scriptsize] at (3,0.18) {\itshape densité de $\chi^2(5)$};
\end{tikzpicture}
\end{center}

\smallskip

\textbf{Statistique du test d'adéquation.} On dispose d'un échantillon de taille $n$ réparti en $r$ catégories d'effectifs observés $O_1, \dots, O_r$. Sous $H_0$ (\og les données suivent telle loi de référence \fg), on calcule les effectifs \emph{attendus} $E_1, \dots, E_r$, puis
\[
  \boxed{\;\chi^2_{\text{obs}} = \sum_{i=1}^{r} \dfrac{(O_i - E_i)^2}{E_i}.\;}
\]
On admet (théorème de Pearson) que sous $H_0$ et pour $n$ grand, $\chi^2_{\text{obs}}$ suit approximativement une loi $\chi^2(r-1)$. La p-value est $P\!\bigl(\chi^2(r-1) \ge \chi^2_{\text{obs}}\bigr)$ : on rejette si elle est petite, c'est-à-dire si l'écart entre observé et attendu est \emph{trop grand} pour être imputé au hasard.
\end{mdframed}

\medskip

\textbf{Énoncé.} Un joueur soupçonne son dé à $6$ faces d'être pipé. Il le lance $n=600$ fois et observe :

\begin{center}
\renewcommand{\arraystretch}{1.3}
\begin{tabular}{|c|c|c|c|c|c|c|}
  \hline
  Face $i$ & $1$ & $2$ & $3$ & $4$ & $5$ & $6$ \\
  \hline
  Effectif observé $O_i$ & $75$ & $90$ & $100$ & $110$ & $100$ & $125$ \\
  \hline
\end{tabular}
\end{center}

\medskip

\textbf{Question 1.} Quelles sont $H_0$ et $H_1$\,? Sous $H_0$, donner les effectifs attendus $E_i$ pour chaque face.

\atrou[3cm]{%
$H_0$ : \og le dé est équilibré \fg, c'est-à-dire $P(\text{face } i) = \frac{1}{6}$ pour tout $i$.\\
$H_1$ : \og le dé est pipé \fg{} (au moins une face n'a pas la probabilité $\frac{1}{6}$).\\
Sous $H_0$, chaque face est attendue $E_i = n \times \frac{1}{6} = \frac{600}{6} = 100$ fois.
}

\textbf{Question 2.} Calculer la statistique $\chi^2_{\text{obs}}$.

\atrou[4.5cm]{%
\begin{align*}
  \chi^2_{\text{obs}} &= \dfrac{(75-100)^2}{100} + \dfrac{(90-100)^2}{100} + \dfrac{(100-100)^2}{100} \\
  &\quad + \dfrac{(110-100)^2}{100} + \dfrac{(100-100)^2}{100} + \dfrac{(125-100)^2}{100} \\
  &= \dfrac{625 + 100 + 0 + 100 + 0 + 625}{100} = \dfrac{1450}{100} = 14{,}5.
\end{align*}
}

\textbf{Question 3.} Quel est le nombre de degrés de liberté\,? Comparer $\chi^2_{\text{obs}}$ à l'espérance d'une loi $\chi^2(k)$ sous $H_0$. Cela paraît-il \og normal \fg{} ou \og suspect \fg{}\,?

\atrou[3cm]{%
Il y a $r=6$ catégories, donc $r-1 = 5$ degrés de liberté. Sous $H_0$, $\chi^2_{\text{obs}}$ suit approximativement $\chi^2(5)$, d'espérance $5$ et d'écart-type $\sqrt{10} \approx 3{,}16$. Or on observe $14{,}5$, soit environ $3$ écarts-types au-dessus de la moyenne attendue : c'est \textbf{très suspect}.
}

\textbf{Question 4.} À l'aide de l'extrait de table de la fonction de répartition de $\chi^2(5)$ ci-dessous, en déduire la p-value, puis conclure aux seuils $\alpha = 5\,\%$ et $\alpha = 1\,\%$.

\begin{center}
\renewcommand{\arraystretch}{1.3}
\small
\begin{tabular}{|c|c|c|c|c|c|c|c|c|c|}
  \hline
  $x$ & $9$ & $10$ & $11$ & $12$ & $13$ & $14$ & $14{,}5$ & $15$ & $16$ \\
  \hline
  $F(x) = P\!\bigl(\chi^2(5) \le x\bigr)$ & $0{,}891$ & $0{,}925$ & $0{,}949$ & $0{,}965$ & $0{,}977$ & $0{,}984$ & $0{,}987$ & $0{,}990$ & $0{,}993$ \\
  \hline
\end{tabular}
\end{center}

\atrou[4.5cm]{%
La p-value est la probabilité d'observer un $\chi^2$ \emph{au moins aussi grand} que $14{,}5$ sous $H_0$ :
\[
  \text{p-value} = P\!\bigl(\chi^2(5) \ge 14{,}5\bigr) = 1 - F(14{,}5) = 1 - 0{,}987 = 0{,}013.
\]
\begin{itemize}
  \item Au seuil $\alpha = 5\,\%$ : $0{,}013 < 0{,}05$, on \textbf{rejette} $H_0$. Le dé semble pipé.
  \item Au seuil $\alpha = 1\,\%$ : $0{,}013 > 0{,}01$, on \textbf{ne rejette pas} $H_0$ : les données sont à la limite, on ne peut pas conclure formellement à $1\,\%$.
\end{itemize}
La conclusion dépend du seuil choisi --- ce qui rappelle l'importance de fixer $\alpha$ \emph{avant} l'expérience.
}

\medskip

\textit{Pour aller plus loin.} Le test du $\chi^2$ s'étend (avec quelques modifications) à de nombreux contextes : indépendance entre deux variables (tableau de contingence), homogénéité entre populations, qualité d'ajustement à une loi continue (après regroupement en classes). On pourra y revenir en TD.

% ====================================================================
\section*{4. Dualité intervalle de confiance $\Longleftrightarrow$ test}
% ====================================================================

\textbf{Question 1.} Au CM4, on avait construit l'intervalle de confiance à $95\,\%$ pour $\mu$ dans le problème des sacs de ciment : $IC_{95\%} = [24{,}702\,;\,24{,}898]$. Cet intervalle contient-il $\mu_0 = 25$\,?

\atrou[2cm]{%
\textbf{Non.} La valeur $25$ est en dehors. C'est exactement la même conclusion que le test : la machine est déréglée.
}

\textbf{Question 2.} Montrer que les deux énoncés suivants sont \textbf{équivalents} :
\begin{itemize}
  \item[(a)] On rejette $H_0 : \mu = \mu_0$ au seuil $\alpha$.
  \item[(b)] $\mu_0 \notin IC_{1-\alpha}$ (l'intervalle de confiance à $1-\alpha$).
\end{itemize}

\atrou[5cm]{%
On a rejet $\Leftrightarrow$ $|Z_{\text{obs}}| > t_\alpha$ $\Leftrightarrow$
$\left|\Xbar - \mu_0\right| > t_\alpha\,\dfrac{\sigma}{\sqrt{n}}$ $\Leftrightarrow$
$\mu_0 \notin \left[\Xbar - t_\alpha\,\dfrac{\sigma}{\sqrt{n}}\,;\,\Xbar + t_\alpha\,\dfrac{\sigma}{\sqrt{n}}\right]$
$\Leftrightarrow$ $\mu_0 \notin IC_{1-\alpha}$.

\smallskip

C'est exactement le même calcul, vu sous deux angles.
}

\begin{mdframed}[backgroundcolor=ThemeLight, linecolor=Theme,
  linewidth=1pt, innertopmargin=8pt, innerbottommargin=8pt, nobreak=true]
\textbf{Dualité IC / test.} Au même seuil $\alpha$ :
\[
  \boxed{\;\text{Rejeter } H_0: \mu = \mu_0 \;\Longleftrightarrow\; \mu_0 \notin IC_{1-\alpha}.\;}
\]
L'intervalle de confiance est l'\textbf{ensemble des valeurs $\mu_0$ qu'on \emph{ne rejetterait pas}} au seuil $\alpha$.
\end{mdframed}

\textit{Conséquence pratique.} Si on a déjà calculé un IC à $95\,\%$, on a \emph{gratuitement} un test bilatéral au seuil $5\,\%$ pour n'importe quelle valeur $\mu_0$.

% ====================================================================
\section*{5. Tests bilatéral et unilatéral}
% ====================================================================

\subsection*{5.1. Définition : trois façons de formuler $H_0$}

Jusqu'ici, on a posé $H_0 : \mu = \mu_0$ contre $H_1 : \mu \neq \mu_0$~: \emph{n'importe quel} écart à $\mu_0$ amène à rejeter. Mais selon la question pratique qu'on se pose, on peut vouloir \textbf{cibler le sens de l'écart} qu'on cherche à détecter. L'hypothèse $H_0$ (et son alternative $H_1$) peut alors prendre trois formes :

\begin{mdframed}[backgroundcolor=ThemeLight, linecolor=Theme,
  linewidth=1pt, innertopmargin=8pt, innerbottommargin=8pt]
\begin{itemize}\setlength{\itemsep}{0.3em}
  \item \textbf{Test bilatéral} : \quad $H_0 : \mu = \mu_0$ \;\;contre\;\; $H_1 : \mu \neq \mu_0$. \\
  On suspecte un écart \emph{dans n'importe quel sens}.
  \item \textbf{Test unilatéral à gauche} : \quad $H_0 : \mu \ge \mu_0$ \;\;contre\;\; $H_1 : \mu < \mu_0$. \\
  On veut seulement détecter un écart \emph{vers le bas} : $H_0$ couvre toutes les valeurs \og acceptables \fg{} ($\mu_0$ et au-dessus).
  \item \textbf{Test unilatéral à droite} : \quad $H_0 : \mu \le \mu_0$ \;\;contre\;\; $H_1 : \mu > \mu_0$. \\
  On veut seulement détecter un écart \emph{vers le haut}.
\end{itemize}
\end{mdframed}

\smallskip

\textit{Remarque.} Pour les calculs, le cas \og pire \fg{} (le plus défavorable au rejet) sous une $H_0$ unilatérale est toujours $\mu = \mu_0$ (la frontière). On garde donc la même statistique $Z = (\Xbar - \mu_0)/(\sigma/\sqrt n) \approx \mathcal{N}(0,1)$ : seule la zone de rejet change.

\smallskip

\textbf{Le contexte dicte le choix.} Le sens de $H_0$ se déduit \emph{toujours} de la situation à analyser, jamais des données observées. Trois exemples :

\begin{itemize}\setlength{\itemsep}{0.2em}
  \item \textbf{Sacs de ciment (CM4).} La machine peut se dérégler \emph{vers le haut} (perte de matière première) ou \emph{vers le bas} (sacs sous-remplis, fraude). On veut détecter \emph{tout} dérèglement~:
  $H_0 : \mu = 25$. \textbf{Bilatéral.}
  \item \textbf{Ampoules du vendeur (CM4).} Le vendeur annonce $\mu_0 = 1000$~h. Si la durée réelle est \emph{supérieure}, c'est tout bénéfice pour l'acheteur, qui n'a aucune raison de se plaindre~: cette situation reste compatible avec $H_0$. On veut seulement attraper le dérapage \emph{vers le bas}~:
  $H_0 : \mu \ge 1000$. \textbf{Unilatéral à gauche.}
  \item \textbf{Nouveau béton.} Un industriel propose un béton qu'il prétend \emph{plus résistant} que le standard à $\mu_0 = 30$~MPa. Le client cherche à valider que le nouveau produit est vraiment meilleur~:
  $H_0 : \mu \le 30$. \textbf{Unilatéral à droite.}
\end{itemize}

\smallskip

\textit{Important.} Le choix \og bilatéral / unilatéral \fg{} se fait \textbf{avant} de regarder les données, jamais pour faire arranger la p-value.

\subsection*{5.2. Conséquence sur la zone de rejet et la p-value}

L'idée est la même dans les trois cas : on rejette $H_0$ si $Z_{\text{obs}}$ tombe \og loin \fg{} dans le(s) sens prévu(s) par $H_1$. La différence est où l'on place les $\alpha$ \% de probabilité \og extrême \fg{} sous $H_0$ :

\begin{center}
\begin{tikzpicture}[scale=0.55]
  % bilateral
  \begin{scope}[xshift=-7cm]
    \draw[thick, Theme, domain=-3.4:3.4, smooth, samples=100]
      plot (\x, {2.2*exp(-\x*\x/2)});
    \fill[Accent, opacity=0.6, domain=-3.4:-1.96, smooth, samples=40]
      plot (\x, {2.2*exp(-\x*\x/2)}) -- (-1.96,0) -- (-3.4,0) -- cycle;
    \fill[Accent, opacity=0.6, domain=1.96:3.4, smooth, samples=40]
      plot (\x, {2.2*exp(-\x*\x/2)}) -- (3.4,0) -- (1.96,0) -- cycle;
    \draw[thick,->] (-4,0) -- (4,0);
    \node[below] at (-1.96,-0.05) {\scriptsize $-1{,}96$};
    \node[below] at (1.96,-0.05) {\scriptsize $1{,}96$};
    \node at (-2.6,1.1) {\scriptsize $\alpha/2$};
    \node at (2.6,1.1) {\scriptsize $\alpha/2$};
    \node at (0,-1.4) {\small\bfseries Bilatéral};
  \end{scope}
  % unilateral gauche
  \begin{scope}[xshift=0cm]
    \draw[thick, Theme, domain=-3.4:3.4, smooth, samples=100]
      plot (\x, {2.2*exp(-\x*\x/2)});
    \fill[Accent, opacity=0.6, domain=-3.4:-1.645, smooth, samples=40]
      plot (\x, {2.2*exp(-\x*\x/2)}) -- (-1.645,0) -- (-3.4,0) -- cycle;
    \draw[thick,->] (-4,0) -- (4,0);
    \node[below] at (-1.645,-0.05) {\scriptsize $-1{,}645$};
    \node at (-2.4,1.1) {\scriptsize $\alpha$};
    \node at (0,-1.4) {\small\bfseries Unilatéral à gauche};
  \end{scope}
  % unilateral droite
  \begin{scope}[xshift=7cm]
    \draw[thick, Theme, domain=-3.4:3.4, smooth, samples=100]
      plot (\x, {2.2*exp(-\x*\x/2)});
    \fill[Accent, opacity=0.6, domain=1.645:3.4, smooth, samples=40]
      plot (\x, {2.2*exp(-\x*\x/2)}) -- (3.4,0) -- (1.645,0) -- cycle;
    \draw[thick,->] (-4,0) -- (4,0);
    \node[below] at (1.645,-0.05) {\scriptsize $1{,}645$};
    \node at (2.4,1.1) {\scriptsize $\alpha$};
    \node at (0,-1.4) {\small\bfseries Unilatéral à droite};
  \end{scope}
\end{tikzpicture}
\end{center}

\medskip

\begin{mdframed}[backgroundcolor=ThemeLight, linecolor=Theme,
  linewidth=1pt, innertopmargin=8pt, innerbottommargin=8pt, nobreak=true]
\textbf{À retenir --- valeurs critiques au seuil $\alpha = 5\,\%$.}
\begin{itemize}\setlength{\itemsep}{0.2em}
  \item \textbf{Bilatéral} ($H_0 : \mu = \mu_0$)~: $\alpha$ se répartit en deux moitiés. \\
  Rejet si $|Z_{\text{obs}}| > 1{,}96$. \quad p-value $= 2\bigl(1-\Phi(|Z_{\text{obs}}|)\bigr)$.
  \item \textbf{Unilatéral à gauche} ($H_0 : \mu \ge \mu_0$)~: tout $\alpha$ dans la queue de gauche. \\
  Rejet si $Z_{\text{obs}} < -1{,}645$. \quad p-value $= \Phi(Z_{\text{obs}})$.
  \item \textbf{Unilatéral à droite} ($H_0 : \mu \le \mu_0$)~: tout $\alpha$ dans la queue de droite. \\
  Rejet si $Z_{\text{obs}} > 1{,}645$. \quad p-value $= 1-\Phi(Z_{\text{obs}})$.
\end{itemize}
\end{mdframed}

\medskip

\textbf{Question 1 (revisite des ampoules, CM4).} Au CM4, on avait $n=50$, $\sigma=100$, $\Xbar=950$, soit
$Z_{\text{obs}} = (950-1000)/(100/\sqrt{50}) = -3{,}53$. Justifier le choix d'un test \emph{unilatéral à gauche}, calculer la p-value associée, et conclure.

\atrou[3.5cm]{%
Le vendeur tire avantage d'une durée \emph{annoncée trop longue} : on ne s'inquiète que du sens \og durent moins \fg, soit $H_0 : \mu \ge 1000$.
\[
  \text{p-value} = \Phi(-3{,}53) \approx 2{,}1\cdot 10^{-4}.
\]
$2{,}1\cdot 10^{-4} \ll 0{,}05$ : on \textbf{rejette} $H_0$. Le vendeur ment.
}

\textbf{Question 2.} Si on avait choisi (à tort) un test bilatéral, quelle aurait été la p-value\,?

\atrou[2.2cm]{%
$\text{p-value}_{\text{bilat}} = 2 \cdot \Phi(-3{,}53) \approx 4{,}2\cdot 10^{-4}$.

\smallskip

\textbf{Deux fois plus grande}, mais on rejette toujours. C'est une règle générale : en passant de bilatéral à unilatéral (du bon côté), la p-value est divisée par deux.
}

\subsection*{5.3. Application : durée de vie d'une élingue de levage}

\textbf{Un mot sur la loi exponentielle.} Pour modéliser une \og durée de vie \fg{} (durée avant panne d'un composant, durée d'un appel, intervalle entre deux clients dans une file d'attente), on utilise très souvent la \textbf{loi exponentielle} $\mathcal{E}(\lambda)$, de densité $f(x) = \lambda e^{-\lambda x}$ pour $x \ge 0$. Ses caractéristiques :
\begin{itemize}\setlength{\itemsep}{0.15em}
  \item $\E(X) = 1/\lambda$ et $\Var(X) = 1/\lambda^2$, donc $\boxed{\sigma = \mu = 1/\lambda}$ : \emph{l'écart-type est égal à la moyenne}. C'est très pratique pour les tests : connaître $\mu$ suffit à connaître $\sigma$.
  \item \textbf{Propriété sans mémoire} : $P(X > s+t \mid X > s) = P(X > t)$. Le composant \og oublie \fg{} son passé : qu'il ait déjà fonctionné $1$~h ou $1000$~h, sa durée de vie résiduelle a la même loi.
  \item Probabilité de survie : $P(X > t) = e^{-\lambda t}$.
\end{itemize}

\smallskip

\textbf{Cadre.} Une entreprise de BTP commande des \textbf{élingues d'acier} pour ses chantiers de levage. Le fabricant annonce une durée de vie moyenne avant rupture de $\mu_0 = 5000$ heures de service. On modélise la durée de vie d'une élingue par une loi $\mathcal{E}(\lambda)$.

Sur un échantillon de $n=50$ élingues testées jusqu'à la rupture, on observe une durée moyenne empirique $\Xbar = 4200$~heures.

\medskip

\textbf{Question 1.} Justifier que le test à mener est \emph{unilatéral} et formuler $H_0$ et $H_1$.

\atrou[2.5cm]{%
Le fabricant aurait intérêt à \emph{surestimer} la durée de vie (argument commercial). Du côté du client, on ne s'inquiète que du dérapage vers le bas :
\[ H_0 : \mu \ge 5000 \quad\text{contre}\quad H_1 : \mu < 5000. \]
\textbf{Unilatéral à gauche.}
}

\textbf{Question 2.} En utilisant le TCL et la propriété $\sigma = \mu$ de la loi exponentielle, construire la statistique de test et calculer $Z_{\text{obs}}$.

\atrou[4cm]{%
Sous $H_0$, $\mu = 5000$ donc $\sigma = 5000$. Par le TCL :
\[
  \Xbar \approx \mathcal{N}\!\left(5000,\;\dfrac{5000^2}{50}\right),
  \qquad
  Z = \dfrac{\Xbar - 5000}{5000/\sqrt{50}} \approx \mathcal{N}(0,1).
\]
\[
  Z_{\text{obs}} = \dfrac{4200 - 5000}{5000/\sqrt{50}} = \dfrac{-800}{707{,}1} \approx -1{,}13.
\]
}

\textbf{Question 3.} Calculer la p-value et conclure au seuil $5\,\%$.

\atrou[3cm]{%
Test unilatéral à gauche :
\[
  \text{p-value} = P(Z \le -1{,}13 \mid H_0) = \Phi(-1{,}13) = 1 - \Phi(1{,}13) \approx 0{,}129.
\]
$0{,}129 > 0{,}05$ : on \textbf{ne rejette pas} $H_0$. L'écart de $800$~h observé est compatible avec une simple fluctuation d'échantillonnage.
}

\medskip

\textit{Moralité --- attention au piège.} \og Ne pas rejeter $H_0$ \fg{} ne veut \emph{pas} dire \og prouver que $H_0$ est vraie \fg. Cela veut dire : \og les données ne suffisent pas à prouver le contraire \fg. Pour trancher, soit on augmente la taille de l'échantillon, soit on accepte de rester sur l'annonce du fabricant par défaut. Si l'enjeu de sécurité est lourd (rupture d'une élingue $\Longrightarrow$ chute de charge), on peut préférer un seuil $\alpha$ \emph{plus grand} ($10\,\%$ par exemple) pour rejeter plus facilement.

% ====================================================================
\section*{6. Deux types d'erreurs}
% ====================================================================

Un test n'est \textbf{jamais} certain : il peut se tromper de deux façons.

\begin{center}
\renewcommand{\arraystretch}{1.6}
\begin{tabular}{|m{4cm}|c|c|}
  \hline
  & $H_0$ vraie en réalité & $H_0$ fausse en réalité \\
  \hline
  Décision : \emph{garder} $H_0$ & \checkmark\ bonne décision & \textbf{erreur de type II} ($\beta$) \\
  \hline
  Décision : \emph{rejeter} $H_0$ & \textbf{erreur de type I} ($\alpha$) & \checkmark\ bonne décision \\
  \hline
\end{tabular}
\end{center}

\medskip

\begin{itemize}
  \item \textbf{Erreur de type I} : on rejette $H_0$ alors qu'elle est vraie. Sa probabilité est précisément $\alpha$ (le seuil qu'on fixe).
  \item \textbf{Erreur de type II} : on garde $H_0$ alors qu'elle est fausse. Sa probabilité $\beta$ dépend de la vraie valeur du paramètre (et n'est en général \emph{pas} égale à $1-\alpha$).
\end{itemize}

\medskip

\textbf{Le compromis fondamental.}

\begin{center}
\begin{mdframed}[backgroundcolor=ThemeLight, linecolor=Theme,
  linewidth=1pt, innertopmargin=8pt, innerbottommargin=8pt]
À $n$ fixé : \;diminuer $\alpha$ $\Longrightarrow$ augmenter $\beta$. \;Diminuer les deux à la fois nécessite d'augmenter $n$.
\end{mdframed}
\end{center}

\medskip

\textbf{Analogie au tribunal.} $H_0$ = \og l'accusé est innocent \fg{}.
\begin{itemize}
  \item Type I : condamner un innocent. Très grave $\Longrightarrow$ on prend $\alpha$ \emph{petit}.
  \item Type II : acquitter un coupable. Moins grave dans la philosophie du droit, mais pas négligeable.
\end{itemize}

\medskip

\textbf{Question --- Application industrielle.}
Une machine produit en moyenne $3\,\%$ de boulons défectueux. On prélève chaque jour $n=200$ boulons. On décide la règle \og \textbf{rejeter le lot et arrêter la machine si plus de $10$ défectueux} \fg.
\begin{enumerate}
  \item Sous $H_0$ (la machine est bien réglée, $p = 0{,}03$), quelle loi suit le nombre $X$ de défauts dans l'échantillon\,? Donner son espérance et son écart-type.
  \item Par approximation normale (TCL), calculer le risque de type I, c'est-à-dire la probabilité d'arrêter la machine alors qu'elle est bien réglée.
\end{enumerate}

\atrou[6cm]{%
\textbf{1.} $X \sim \mathcal{B}(200,\,0{,}03)$, donc
\[
  \E(X) = np = 6,\qquad \sigma(X) = \sqrt{np(1-p)} = \sqrt{200\cdot 0{,}03\cdot 0{,}97} \approx 2{,}41.
\]

\textbf{2.} Par le TCL, sous $H_0$, $X \approx \mathcal{N}(6,\,2{,}41^2)$.
\[
  P(X > 10 \mid H_0) = P\!\left(Z > \dfrac{10-6}{2{,}41}\right) = P(Z > 1{,}66) \approx 1-\Phi(1{,}66) \approx 0{,}049.
\]
Soit environ $5\,\%$ : c'est notre seuil $\alpha$. Une fois tous les vingt jours environ, on arrêtera la machine \emph{à tort}.
}

\medskip

\textit{Et si on durcissait la règle\,?} Avec \og rejet si plus de $14$ défauts \fg{}, $\alpha$ tomberait à environ $0{,}05\,\%$... mais on laisserait passer beaucoup plus de lots vraiment défectueux ($\beta$ augmente). C'est le compromis.

% ====================================================================
\needspace{20\baselineskip}
\section*{7. Bilan}
% ====================================================================

\nopagebreak
\begin{mdframed}[backgroundcolor=ThemeLight, linecolor=Theme,
  linewidth=1.5pt, innertopmargin=10pt, innerbottommargin=10pt, nobreak=true]
\textbf{À retenir.}
\begin{enumerate}
  \item \textbf{Test d'hypothèse} : on confronte une hypothèse $H_0$ (modèle de référence) aux données. On rejette $H_0$ si l'observation est \og trop improbable \fg{} sous $H_0$.

  \item \textbf{p-value} : $P(\text{observation au moins aussi extrême} \mid H_0)$. On rejette si $\text{p-value} < \alpha$.

  \item \textbf{Recette d'un test sur la moyenne ($\sigma$ connu)} : statistique $Z = (\Xbar - \mu_0)/(\sigma/\sqrt{n})$, qui suit $\mathcal{N}(0,1)$ sous $H_0$. Rejet bilatéral si $|Z_{\text{obs}}| > t_\alpha$ ($t_{0{,}05}=1{,}96$) ; rejet unilatéral si $Z_{\text{obs}}$ franchit $\pm u_\alpha$ ($u_{0{,}05}=1{,}645$).

  \item \textbf{Test du khi-deux d'adéquation} : $\chi^2_{\text{obs}} = \sum (O_i - E_i)^2/E_i$, distribué $\chi^2(r-1)$ sous $H_0$. Rejet si $\chi^2_{\text{obs}}$ \og trop grand \fg.

  \item \textbf{Dualité IC $\Leftrightarrow$ test} : on rejette $H_0 : \mu = \mu_0$ au seuil $\alpha$ si et seulement si $\mu_0 \notin IC_{1-\alpha}$.

  \item \textbf{Bilatéral vs unilatéral} : choisir d'après le contexte, \emph{avant} l'expérience.

  \item \textbf{Deux types d'erreurs} : type I ($\alpha$, fixé) et type II ($\beta$, dépend de la vraie valeur). Pour diminuer les deux à la fois : augmenter $n$.
\end{enumerate}
\end{mdframed}

\end{document}
